\begin{helper}
Fix a history cell \(H\), a terminal coordinate set \(U\), a nodup terminal
order \(o\) with \(\operatorname{toFinset}(o)=U\), a blue support label
\(B\subseteq U\) of size \(u_B\), and a fixed red label \(J\).  Suppose the
branch/transcript interface supplies finite branch labels \(\beta\), finite
transcript labels \(\tau=(P_\tau,A_\tau,Z_\tau)\), complete terminal blue
traces \(\operatorname{blueTrace}(\beta)\), decoded red supports
\(R_{\beta,J}\), a realized-cell predicate
\(\mathsf{Realized}(\beta,\tau)\), and a compatibility predicate for
complete terminal assignments.

Assume the terminal product law for \(H\), \(0\le p\le1/2\), and the
fixed-\((B,J)\) geometry supplied by the branch/transcript construction:
for every realized \((\beta,\tau)\),
\[
|R_{\beta,J}|=u_R,\qquad R_{\beta,J}\subseteq U,
\]
\[
P_\tau,A_\tau\subseteq U,\qquad P_\tau\cap A_\tau=\emptyset,
\qquad B\cup R_{\beta,J}\subseteq P_\tau,
\]
\[
Z_\tau\subseteq A_\tau,\qquad
\max(0,a-u_R-u_B)\le |Z_\tau|,
\]
and \(P_\tau,A_\tau\) record the complete terminal blue trace of
\(\beta\).  Compatibility is forward-correct for the residual cylinder:
compatible complete assignments contain \(P_\tau\), avoid
\(A_\tau\setminus Z_\tau\), and agree with the complete blue trace away from
the selected set.  Also, one complete terminal assignment is compatible with
at most one realized pair.

Assume in addition the explicit residual partition data for this fixed
\((B,J)\).  There is a nonnegative residual assignment mass \(\mu(A)\) on
complete terminal assignments \(A\subseteq U\), with
\[
\sum_{A\subseteq U}\mu(A)\le 1,
\]
and, for each realized \((\beta,\tau)\), a finite released block
\(\mathcal R_{\beta,\tau}\subseteq\{A:A\subseteq U\}\) and a block mass
\(\rho_{\beta,\tau}\) such that
\[
\rho_{\beta,\tau}
  =\sum_{A\in\mathcal R_{\beta,\tau}}\mu(A).
\]
Released blocks for distinct realized pairs are disjoint.  Finally, every
realized pair satisfies the residual domination
\[
p^{|P_\tau|-u_R-u_B}(1-p)^{|A_\tau|-|Z_\tau|}
  \le \rho_{\beta,\tau}.
\]

Then there is a finite residual leaf set \(\Lambda_{B,J}\), a residual leaf
map \(\iota(\beta,\tau)\in\Lambda_{B,J}\), and nonnegative leaf masses
\(\nu\) such that
\[
\sum_{\lambda\in\Lambda_{B,J}}\nu(\lambda)\le1,
\]
the leaf map is injective on realized pairs, and every realized pair
satisfies
\[
p^{|P_\tau|-u_R-u_B}(1-p)^{|A_\tau|-|Z_\tau|}
  \le \nu(\iota(\beta,\tau)).
\]
\end{helper}
\begin{proof}
Apply the finite released-cylinder packing lemma
\(\noderef{FixedSetHistoryCellRedReleasedCylinderPacking}\) to the complete
terminal assignment space \(2^U\), the mass \(\mu\), and the released blocks
\(\mathcal R_{\beta,\tau}\).  The hypotheses of that node are exactly the
normalization data listed above: \(\mu\) is nonnegative, the total
\(\mu\)-mass of \(2^U\) is at most \(1\), each released block is a subset of
\(2^U\), the block mass is the \(\mu\)-sum over that block, and distinct
realized blocks are disjoint.  Hence
\[
\sum_{\beta}
  \sum_{\tau\in\mathcal P:\mathsf{Realized}(\beta,\tau)}
    \rho_{\beta,\tau}
\le 1 .
\]

Define the residual leaf set to be the finite product of the branch-label
type and the transcript-label type.  Send \((\beta,\tau)\) to its own pair
\[
\iota(\beta,\tau)=(\beta,\tau).
\]
Give a leaf \((\beta,\tau)\) mass \(\rho_{\beta,\tau}\) when
\(\tau\) is an allowed transcript and \((\beta,\tau)\) is realized, and give
it mass \(0\) otherwise.  Nonnegativity follows because realized
\(\rho_{\beta,\tau}\) is a sum of nonnegative \(\mu\)-masses, and all
non-realized leaves receive mass \(0\).

The total leaf mass is exactly the displayed double sum over realized
allowed pairs, since all other leaves have mass \(0\).  Therefore the
released-cylinder packing bound gives
\[
\sum_{\lambda\in\Lambda_{B,J}}\nu(\lambda)\le1 .
\]
If two realized pairs have the same leaf, then the ordered pairs
\((\beta,\tau)\) are equal, so both the branch label and transcript are
equal.  Thus the leaf map is injective on realized pairs.

For a realized pair, the mass of its leaf is
\(\nu(\iota(\beta,\tau))=\rho_{\beta,\tau}\).  The assumed residual
domination of \(\rho_{\beta,\tau}\) is therefore exactly
\[
p^{|P_\tau|-u_R-u_B}(1-p)^{|A_\tau|-|Z_\tau|}
  \le \nu(\iota(\beta,\tau)).
\]
This proves the residual leaf certificate.  The selected set \(Z_\tau\) is
used only through the residual domination hypothesis, so selected blue
absences and selected red absences are handled uniformly; no red-only
selected-set assumption is introduced in this packaging step.
\end{proof}
